
Or @@M0@@ est de
rang @@M1@@, donc @@M2@@, d'où le résultat.

\subsection*{4.2.}

Soit @@M3@@ une immersion régulière. Soit @@M4@@ un entier @@M5@@, et @@M6@@ la section définie par l'homomorphisme : @@M7@@, passage au quotient par le sous-faisceau engendré par les @@M8@@ dernières coordonnées. Si @@M9@@ est défini par la @@M10@@-Algèbre graduée @@M11@@, l'idéal @@M12@@ définissant la section @@M13@@ est l'idéal engendré par les sections @@M14@@, donc @@M15@@ est isomorphe à @@M16@@, donc l'idéal correspondant de @@M17@@ à @@M18@@.

Comme la section @@M19@@ est évidemment une immersion régulière, il en est de même pour @@M20@@, et on peut écrire la suite exacte des faisceaux conormaux correspondante (1.7). On obtient ainsi l'égalité suivante dans @@M21@@, en notant @@M22@@ le faisceau conormal de @@M23@@, et @@M24@@ le faisceau conormal de @@M25@@ :
@@M26@@
On voit donc que pour toute immersion régulière @@M27@@, on peut, en la composant avec une section : @@M28@@, ou bien : @@M29@@, obtenir une immersion régulière satisfaisant la propriété du lemme 4.1.

\textbf{Théorème 4.3.} \uline{Soient} @@M30@@ \uline{une immersion fermée régulière}, @@M31@@ \uline{la classe dans} @@M32@@ \uline{du faisceau conormal de} @@M33@@ \uline{dans} @@M34@@. \uline{On suppose que} @@M35@@ \uline{est quasi-compact et a un faisceau ample}. \uline{Alors}, \uline{avec les notations de} 3.3, \uline{on a pour tout} @@M36@@, \uline{et tout} @@M37@@ :
@@M38@@
@@M39@@ \uline{étant l'élément défini dans} (V 5.3).

i) On suppose d'abord que l'hypothèse de 4.1 est satisfaite ; on va dans ce cas montrer que @@M40@@. Comme @@M41@@, cette
